% !TeX program = xelatex
% ---------------------------------------------------------------------------
% Guarantee or Alarm: Statistical Verification for LLM-Assisted Decision
% Pipelines. A three-layer framework with exact anchors, perturbation
% checks, and runnable evidence.
%
% Compiles with XeLaTeX (MiKTeX / TeX Live) and also with pdfLaTeX
% (the iftex guard below switches font handling automatically).
% Build: xelatex main && bibtex main && xelatex main && xelatex main
% ---------------------------------------------------------------------------
\documentclass[11pt,a4paper]{article}

\usepackage{iftex}
\ifPDFTeX
  \usepackage[T1]{fontenc}
  \usepackage[utf8]{inputenc}
  \usepackage{lmodern}
\else
  \usepackage{fontspec}
  % Defaults to Latin Modern (always available). If the tex-gyre package is
  % installed (MiKTeX Console -> Packages -> tex-gyre), uncomment for Pagella:
  % \setmainfont{TeX Gyre Pagella}
  % \setmonofont[Scale=0.88]{TeX Gyre Cursor}
\fi

\usepackage[margin=2.7cm]{geometry}
\usepackage{amsmath,amssymb,amsthm}
\usepackage{booktabs}
\usepackage{array}
\usepackage{microtype}
\setlength{\emergencystretch}{1.2em}
\usepackage{enumitem}
\usepackage[numbers,sort&compress]{natbib}
\usepackage{xcolor}
\usepackage[colorlinks=true,linkcolor=blue!50!black,citecolor=blue!50!black,urlcolor=blue!50!black]{hyperref}

% Skill names are typeset in monospace but may break after hyphens
% (url-package machinery), so long names fit narrow table columns.
\DeclareUrlCommand\skillname{\urlstyle{tt}\def\UrlBreaks{\do\-}}
\newcolumntype{L}[1]{>{\raggedright\arraybackslash}p{#1}}

\title{Guarantee or Alarm: Statistical Verification\\
for LLM-Assisted Decision Pipelines\\[3pt]
\large A Three-Layer Framework with Exact Anchors,\\
Perturbation Checks, and Runnable Evidence}

\author{Alvaro Moraes\thanks{Independent researcher.
E-mail: \texttt{alvaromoraes.ici@gmail.com}.}}

\date{July 2026}

\begin{document}
\maketitle

\begin{abstract}
Outputs produced with the help of large language models (LLMs) ---
extractions, labels, estimates, screens, simulations, recommendations ---
increasingly feed consequential decisions. Their characteristic failures are
statistical and \emph{silent}: the pipeline runs, the number is plausible,
and the loss is realized downstream, where no reviewer connects it to the
uncalibrated judge, the peeked-at dashboard, or the data-selected
``winner'' that caused it. We argue that the appropriate response is the
verification-and-validation (V\&V) discipline of computational science,
transplanted to LLM-derived outcomes, and we present a working
implementation. The framework has three layers: \emph{verification of
tooling} (every statistical procedure an agent carries must reproduce an
exactly known ground truth --- the analogue of the method of exact
solutions); \emph{operating-characteristic verification} (demonstrated
coverage, false-discovery rate, and false-alarm rate on synthetic ground
truth, with validation proper --- gold-set audits on the live workload ---
prescribed but distinguished); and \emph{runtime V\&V of outputs} (every
deliverable carries a stated guarantee where one exists, feeding a named
loss, or a diagnostic alarm that halts and escalates). The implementation is a library of eleven agent skills ---
modular, versioned playbooks with dependency-free executable scripts ---
covering calibrated acceptance gates for LLM verifiers, anytime-valid
monitoring, prediction-powered inference from model labels, optimal
experimental design, rare-event and risk-measure estimation, selective and
causal inference, simulation output analysis, and cost-tiered model
orchestration. Twelve numerical demonstrations --- one per use case, plus perturbation
companions and a three-skill composition chain --- contrast naive and
V\&V-informed analyses on simulated data with known ground truth; the gaps are large
(nominal-$90\%$ intervals with $0$--$7\%$ realized coverage repaired to
near nominal; a $47\%$ annual false-alarm rate reduced to $2\%$;
$3\times$ overstated causal effects corrected), and perturbation checks
show what breaks, and which alarms fire, when the assumptions behind the
guarantees are violated. We also describe how
the library was built --- a curation protocol that harvests the implicit
curriculum of methodological schools in the stochastics literature under
adversarial admission control --- and report its economics. All skills are
reusable by any agent framework that can read markdown and run Python.
\end{abstract}

\section{Introduction}\label{sec:intro}

LLM-assisted work fails statistically long before it fails syntactically.
An agent that writes flawless code will still happily report an
uncertainty-free point estimate, treat its own model-generated labels as
ground truth, select a model on the same data it evaluates it on, stop an
experiment the first time a dashboard looks significant, or average its way
past a heavy tail. None of these failures raises an exception. Each
produces a plausible number. And increasingly, those numbers feed decisions
that matter: what to ship, whom to treat, where to provision, what to
certify.

Computational science confronted the same problem decades ago --- complex
software producing plausible numbers for consequential decisions --- and
developed a discipline for it: \emph{verification and validation} (V\&V)
\citep{roache1998,oberkampf2010}. Verification asks whether the equations
are solved right (code verification against exact solutions; solution
verification by a-posteriori error estimation); validation asks whether the
right equations are solved for the intended use. The discipline's central
insight is procedural: correctness claims about computed quantities are not
established by inspection or authority but by \emph{falsifiable checks
attached to the artifact itself}.

This paper transplants that discipline to LLM-derived outcomes and reports
a working implementation. The transplant is direct
(Section~\ref{sec:framework}):

\begin{itemize}[leftmargin=1.6em]
\item \textbf{Verification of tooling.} Every statistical procedure the
agent carries ships with a \emph{verification anchor}: an exactly known
ground truth (closed form, exact finite-sample distribution, analytically
tractable special case) that the shipped code must reproduce --- the
analogue of the method of exact or manufactured solutions. The anchor
certifies the \emph{reference implementation}; Section~\ref{sec:framework}
is explicit about what it does not certify.
\item \textbf{Operating-characteristic verification.} Every procedure
demonstrates its operating characteristics --- coverage, false-discovery
rate, false-alarm rate, cost per verified artifact --- on synthetic data
with known ground truth, with the naive analysis run alongside as control
and with perturbation checks that violate the procedure's assumptions.
This is deliberately \emph{not} called validation: validation proper ---
comparison against audited data from the intended workload --- is
prescribed by the framework as a continuing gold-set audit and is not
claimed by this paper.
\item \textbf{Runtime V\&V of outputs.} Every number delivered to a
decision carries either a \emph{guarantee} (a stated bound, of a stated
kind, feeding a named loss) or an \emph{alarm} (a diagnostic that halts
the pipeline and escalates to a human). Silent invalidity is the failure
mode being engineered away.
\end{itemize}

The implementation vehicle is the \emph{agent skill}
\citep{anthropic2025skills}: a modular, versioned package of procedural
knowledge --- a markdown playbook plus dependency-free executable scripts
--- that an agent loads on demand. Skills are the right carrier for V\&V
because they are inspectable, diffable, and testable: a skill is a claim
that a procedure is correct, and the anchor bundled with it is the means of
falsifying that claim. The library presented here comprises eleven skills
spanning the statistical needs of LLM-assisted decision pipelines, from
calibrated acceptance gates for LLM verifiers to risk-measure estimation
(Section~\ref{sec:toolkit}), each grounded in an identifiable literature
and each verified against an exact anchor
(Section~\ref{sec:verify}).

Two further contributions support the main one. First, twelve numerical
demonstrations --- one per use case plus perturbation companions and a
three-skill composition chain, standard-library Python, fixed seeds,
simulated ground truth --- measure the gap between naive and V\&V-informed
analysis and show, under perturbation, which alarms fire when assumptions
break (Section~\ref{sec:toolkit}, Tables~\ref{tab:numbers}
and~\ref{tab:perturb}); the gaps are not subtle. Second, we describe how the library was built: a curation
protocol, \emph{curriculum harvesting}, that grows the library from
methodological schools in the stochastics and uncertainty-quantification
(UQ) literature under adversarial admission control
(Section~\ref{sec:harvest}), together with the cost-tiered orchestration
that makes the construction itself economical
(Section~\ref{sec:orchestration}). The protocol is secondary to the V\&V
thesis but is, we believe, the natural supply chain for it: the stochastics
literature is unusual in how often it provides exactly the checkable
special cases that anchors require.

\paragraph{Scope and provenance.} The concrete host system is Anthropic's
Claude agent stack (Claude Code and its skill format), with a
multi-provider model panel accessed through the OpenRouter API for cheap
bulk work; nothing in the framework is specific to that stack.
Quantitative claims about the library's construction are drawn from the
project's append-only run logs; the numerical demonstrations are fully
reproducible from the supplementary scripts. Both should be read as a
\emph{case report with runnable evidence}, not as a controlled evaluation.

\section{Background}\label{sec:background}

\subsection{V\&V in computational science}

The V\&V tradition \citep{roache1998,oberkampf2010,sargent2013}
distinguishes: \emph{code verification} --- demonstrating that the
implementation solves its governing equations correctly, canonically by
manufacturing or selecting problems with exact solutions and measuring
convergence to them; \emph{solution verification} --- estimating the
numerical error of a particular computed answer, typically a posteriori;
and \emph{validation} --- demonstrating, by comparison against data from
the intended domain of use, that the model is adequate for the decision it
informs. The tradition is uncompromising on one point: these are distinct
activities, and evidence for one is not evidence for another. A code can be
verified and the model invalid; a model can be well-validated in one regime
and worthless in the next.

LLM-assisted analysis needs exactly this decomposition, for exactly the
original reason: the artifact is too complex for correctness-by-inspection,
and its consumers are too far downstream to catch errors informally. What
changes in the transplant is the object: not a PDE solver but a
\emph{pipeline of statistical procedures}, some executed by code the agent
writes, some by judgments the agent (or a second model) makes.

One element of the transplant must be stated with care, because the
classical vocabulary is strict. In classical V\&V, \emph{validation}
requires comparison against data from the intended operating environment;
checks against a synthetic truth the analyst constructed --- however
adversarial --- are forms of \emph{verification}. This paper's second
layer is therefore named operating-characteristic \emph{verification},
not validation; what this paper's evidence establishes is that the tooling
is correct and that the procedures deliver their advertised error rates on
ground-truth-known problems, including under assumption violations. The
framework \emph{prescribes} validation --- continuing gold-set audits on
the live workload, in the spirit of \citet{sargent2013} --- and
Section~\ref{sec:limitations} is explicit that no such audit is reported
here. The transplant claims the discipline's structure, not credit for
its final stage.

\subsection{Agent skills as procedural memory}

A skill in the sense used here is a directory with a short frontmatter
description (always in the agent's context) and a body loaded only on
trigger, optionally with reference notes and scripts
\citep{anthropic2025skills}. Functionally, skills are \emph{procedural
memory}: they change what the agent does, not what it knows, which
distinguishes them from retrieval-augmented generation (facts at answer
time) \citep{lewis2020} and from fine-tuning (neither inspectable nor
versionable by the end user). For V\&V purposes their decisive property is
that they are artifacts: a skill can be code-reviewed, version-bumped,
diffed --- and shipped with its own falsification test.

\subsection{Why statistical failure is the binding failure mode}

For a deployment playbook, the cost of a subtly wrong skill is a failed
build: loud, fast, cheap. For a statistical procedure the failure mode is
inverted: wrong procedures produce plausible numbers silently, and the loss
is realized far downstream. This asymmetry is why the library's admission
standard is stricter for statistical skills than any marketplace's, and why
the runtime layer insists on guarantee-or-alarm: an agent operating
autonomously has no statistician looking over its shoulder, so the
statistician's scruples must be installed as procedure.

The library is organized under a base meta-skill
(\skillname{statistical-grounding}, shipped with the library as its
twelfth package) encoding an eight-rule uncertainty-and-decision contract --- every estimate carries uncertainty;
every decision names its loss; every selection procedure accounts for the
selection; and so on. The individual skills of
Section~\ref{sec:toolkit} are the machinery that makes specific rules
executable.

\section{The V\&V framework for LLM-derived outcomes}\label{sec:framework}

\subsection{Layer 1: verification of tooling (anchors)}

Every skill that ships numeric code must reproduce a ground truth exactly
known by other means, \emph{on the same code path a user will exercise}.
This is code verification in the classical sense, and the anchors play the
role of exact solutions: the expected information gain of a linear-Gaussian
design has a closed form; the Knothe--Rosenblatt transport map of a
Gaussian is its Cholesky factor; split-conformal coverage at calibration
size $n$ is exactly $\lceil (n{+}1)(1{-}\alpha)\rceil/(n{+}1)$ under
exchangeability. Section~\ref{sec:verify} details the anchors and reports a
bug caught by one: approximate checks pass code that exact checks kill.

\subsection{Layer 2: operating-characteristic verification (and where
validation proper begins)}

An anchor certifies the code; it does not certify that the procedure
delivers its advertised error rates on workload-like data, under
workload-like misuse. Layer 2 requires each procedure to demonstrate its
operating characteristics --- realized coverage of nominal intervals,
realized FDR against the budget, realized false-alarm rate under
continuous monitoring, realized cost per correct artifact --- on simulated
data with known ground truth, with the \emph{naive} analysis run alongside
as the control, and with \emph{perturbation checks} that deliberately
violate the procedure's assumptions to show what breaks and which alarm
fires. The demonstrations of Section~\ref{sec:toolkit} are this layer,
executed. Validation proper begins where synthetic truth ends: on live
workloads the framework prescribes continuing gold-set audits ---
periodically label a random sample and compare the pipeline's claims
against it --- and until such an audit exists for a given workload, the
framework treats the pipeline's guarantees as conditional on assumptions
no one has checked against reality.

\paragraph{What counts as a guarantee.} Not all guarantees are equal, and
conflating them is itself a silent failure. The library distinguishes:
\emph{exact finite-sample, distribution-free} (split conformal under
exchangeability; e-processes and confidence sequences under the martingale
assumption; the Clopper--Pearson bound of the verifier gate);
\emph{finite-sample, model-based} (linear-Gaussian design posteriors);
\emph{asymptotic or semiparametric} (DML/AIPW intervals, batch-means
$t$-intervals); and \emph{diagnostic-only} (importance-sampling efficiency
gates, transport-map adequacy, cascade disagreement), which never ship as
guarantees at all --- they only fire alarms. Every skill's playbook states
which kind it delivers, and a deliverable's guarantee class travels with
it to the decision.

\subsection{Layer 3: runtime V\&V of outputs (guarantee or alarm)}

The delivery rule: a number reaching a decision-maker carries a stated
guarantee --- of a stated kind, under stated assumptions, \emph{feeding a
named loss} --- or it does not ship; instead an alarm fires, stating which
assumption failed or which diagnostic tripped, and the case escalates to a
human. The loss clause matters: a defect-rate bound or a coverage
statement is an input to a decision, not a decision, and the library's
base contract (``every decision names its loss'') requires the guarantee
to arrive attached to the cost model it feeds --- defect cost versus
abstention cost at a gate, false-alarm budget versus detection delay on a
monitor, experiment cost versus information at a design. Solution
verification has a direct analogue here: error indicators drive
\emph{escalation}, as in the cost-tiered orchestration of
Section~\ref{sec:orchestration}, where disagreement between cheap and
mid-tier models is itself the a-posteriori error signal that routes work
to a stronger model.

\subsection{Scope: what the framework does not verify}
\label{sec:scope}

Two failure surfaces are explicitly outside what the anchors certify, and
naming them is part of the framework's honesty. First,
\emph{agent-execution fidelity}: the anchor verifies the reference script,
but the agent decides when to invoke it, constructs its arguments, and
interprets its output --- and an agent can select the right script for the
wrong problem, pass malformed inputs, or misstate a valid result in prose.
The library's mitigations are structural (skills forbid re-implementation;
guarantees are computed inside the verified script, so the glue can
suppress a deliverable but not forge one; script boundaries carry input
validation; alarms are emitted as machine-readable artifacts --- exit
codes and structured records --- so a harness, not the agent's prose, is
their intended consumer) but partial: an agent relaying results in prose
can still misstate or omit an alarm to a human reader, so deployments
should route alarm artifacts around the agent, and verification of the
agent's invocation behavior remains an open problem. The paper's claims
should be read as covering the statistical tooling of LLM-assisted
decisions, not the agent's reasoning.
Second, \emph{composition}: statistical validity does not compose ---
upstream imputation, selection, or labeling can invalidate the
exchangeability or ignorability a downstream skill needs. The library's
current discipline is a worked chained demonstration
(Section~\ref{sec:toolkit}, ex10) in which a monitoring link detects
upstream drift, plus per-skill precondition statements; contract-level
composition semantics are named as the frontier, not claimed.

The three layers are ordered by scope --- the first certifies the tools,
the second the procedures, the third each delivered number --- and they
fail independently, which is why all three are needed. The next section
shows the framework at work.

\section{The toolkit in action}\label{sec:toolkit}

A skill earns its place where statistical malpractice is silent: the code
runs, the number is plausible, and the loss lands downstream. The common
thread of the vignettes below is the conversion of a silent invalidity into
a guarantee or an alarm. Table~\ref{tab:usecases} maps the toolkit;
Table~\ref{tab:skills} lists each skill's literature grounding. Some uses
are deployed inside the library's own construction pipeline (and marked
so); the rest are prospective and stated as such.

Every vignette is accompanied by a numerical demonstration: a companion
script (Python standard library only, fixed seed, seconds of runtime on a
laptop) that simulates a dataset with \emph{known ground truth} and runs
the naive and the V\&V-informed analysis side by side. This is Layer~2 of
the framework, executed. Table~\ref{tab:numbers} collects the headline
numbers; each vignette quotes its own; the scripts accompany the paper as
supplementary material (\texttt{examples/}, ex1--ex10, including the
gate's shipped drift monitor ex1b and the nested-risk companion ex5b).

What the demonstrations do and do not establish must be said before the
table, not after it. Each data-generating process was \emph{chosen} so
that ground truth is exact and the contrast is legible; each ``naive''
column is a documented default of current LLM pipelines, but it is also a
textbook anti-pattern meeting its canonical correction on the correction's
home turf. The demonstrations are therefore \emph{reproducible
counterexamples} --- mechanism proofs that these failure modes are generic
and mechanized, and regression tests for the library --- not effect-size
estimates for any real workload, and not an evaluation of the framework
against a no-framework baseline. What separates the columns is correct
versus incorrect statistical practice; the framework's own contribution is
making the correct column \emph{agent-discoverable and runtime-audited}
--- the methods packaged so an agent finds them at the moment of need,
with guarantees computed inside verified code and known boundaries
alarmed --- and the perturbation checks of Table~\ref{tab:perturb}, which
violate each demonstration's key assumption and report what breaks and
which alarm fires, are the part of the evidence that speaks to the
framework rather than to the methods it carries.

\begin{table}[t]
\centering\small
\caption{Use-case map: the silent failure each skill replaces, and what the
practitioner gets instead.}
\label{tab:usecases}
\begin{tabular}{@{}L{3.3cm}L{5.2cm}L{5.2cm}@{}}
\toprule
Skill & Silent failure of naive practice & What the skill delivers \\
\midrule
\skillname{conformal-uq} & LLM judge accepts artifacts at an unknown,
drifting error rate & Finite-sample bound on the accepted-defect rate;
calibrated abstention \\
\skillname{anytime-valid} & Continuously watched dashboards guarantee a
false alarm under peeking & Inference valid at any stopping time; alarm
budget for streams of hypotheses \\
\skillname{prediction-powered-inference} & Model-generated labels treated
as ground truth; bias of unknown sign & Valid intervals at near-model-scale
precision from a small gold set \\
\skillname{bayesian-optimal-design} & Experiment budget spent on
low-information measurements & Designs ranked by expected information gain
per unit cost \\
\skillname{rare-event-is},
\skillname{risk-measure-estimation} & Tail probabilities and CVaR from
sample sizes that have never seen the tail & Feasible certification of
small probabilities; risk functionals with honest error \\
\skillname{selective-inference} & Confidence intervals reported on
data-selected winners & Selection-adjusted intervals; FDR-controlled
shortlists \\
\skillname{causal-ml} & Correlational uplift read as causal; one average
effect for everyone & Confounding-robust, heterogeneous effects; policies
with valid inference \\
\skillname{simulation-output-analysis},
\skillname{measure-transport-inference} & Simulator output averaged
without regard to autocorrelation or transient & Decision-grade error bars
on simulation summaries; fast recalibration \\
\skillname{model-hierarchy} & Frontier-model spend on work a cheap model
verifiably does & Cost per \emph{verified} artifact as the managed
quantity \\
\bottomrule
\end{tabular}
\end{table}

\begin{table}[t]
\centering\small
\caption{The skill library with the primary literature each skill rests
on.}
\label{tab:skills}
\begin{tabular}{@{}L{4.3cm}L{9.9cm}@{}}
\toprule
Skill & Owns / anchored in \\
\midrule
\skillname{model-hierarchy} & Cost-tiered orchestration of LLM workloads as
an MLMC-style telescoping system \citep{giles2008,giles2015,hajiali2016} \\
\skillname{bayesian-optimal-design} & Expected-information-gain estimator
ladder for design of experiments \citep{long2013,beck2018,ryan2016};
sequential OED as a belief MDP \\
\skillname{rare-event-is} & Importance sampling and splitting for small
probabilities; efficiency diagnostics \citep{bucklew2004,rubino2009} \\
\skillname{measure-transport-inference} & Triangular (Knothe--Rosenblatt)
transport maps for Bayesian inference
\citep{elmoselhy2012,marzouk2016,spantini2018} \\
\skillname{conformal-uq} & Split-conformal prediction and calibrated verifier
gates \citep{vovk2005,lei2018,angelopoulos2023gentle} \\
\skillname{selective-inference} & Multiplicity, knockoffs, and
post-selection validity \citep{bh1995,barber2015,candes2018} \\
\skillname{anytime-valid} & Confidence sequences, e-values, online FDR
\citep{howard2021,ramdas2023,waudbysmith2024} \\
\skillname{prediction-powered-inference} & Valid inference from
ML-predicted labels plus a small gold set \citep{angelopoulos2023ppi} \\
\skillname{causal-ml} & Heterogeneous effects, DML, policy learning
\citep{wager2018,chernozhukov2018,athey2021} \\
\skillname{simulation-output-analysis} & Steady-state output CIs;
Rhee--Glynn unbiased multilevel randomization
\citep{asmussen2007,rhee2015} \\
\skillname{risk-measure-estimation} & VaR/CVaR and nested-simulation
estimation; elicitability \citep{rockafellar2000,gneiting2011,giles2019} \\
\bottomrule
\end{tabular}
\end{table}

\begin{table}[t]
\centering\small
\caption{Numerical demonstrations (simulated ground truth, seed 2026;
companion scripts in the supplement). Coverage figures are for nominal
$90\%$ intervals.}
\label{tab:numbers}
\begin{tabular}{@{}L{3.4cm}L{5.1cm}L{5.1cm}@{}}
\toprule
Demonstration & Naive analysis & V\&V-informed analysis \\
\midrule
Verifier gate (ex1) & $8.0\%$ defects among accepted; over the $5\%$
target in $100\%$ of runs & $2.8\%$ defects; over target in $8.3\%$ of runs
(guarantee: $\le 10\%$) \\
Dashboard, null (ex2) & $47.4\%$ chance of $\ge 1$ false alarm per year of
daily peeking at $\alpha=.05$ & $2.2\%$ ($\le 5\%$ guaranteed at any
stopping time) \\
Dashboard, effect (ex2) & fixed-horizon design waits day 365 & confidence
sequence stops at median day 117 \\
Model labels (ex3) & model-only: coverage $0\%$; gold-only: CI width 6.2 pts
& PPI: width 2.9 pts, coverage $89.7\%$ ($4.5\times$ the gold sample) \\
Experiment design (ex4) & random placement: median 37 experiments to target
posterior precision; repeated-center: never & greedy-EIG: 22 experiments \\
Rare event, $p\!\approx\!10^{-6}$ (ex5) & $9.8\times 10^{7}$ samples for
$10\%$ relative error & $5.4\times 10^{2}$ samples (variance reduction
$1.8\times 10^{5}$) \\
Nested risk (ex5b) & inner $M{=}1$: estimates $0.231$ for a true $0.050$ &
budget-aware $M{=}100$: $0.053$, $10\times$ lower RMSE at equal cost \\
Selected winners (ex6) & naive CIs on selected nulls: coverage $0\%$ &
conditional (truncated-normal) CIs: $90.4\%$; BH shortlist at certified
FDR $\le 10\%$ \\
Causal effect (ex7) & difference in means: $6.23$ for a true ATE of $2.00$
& cross-fitted DML/AIPW: $2.00 \pm 0.04$; targeted policy worth $2\times$
treat-all \\
Queue summary (ex8) & iid-style CI: coverage $5.0\%$ (width $0.16$ around a
truth of $9$) & batch means: coverage $83.3\%$ (width $2.7$) \\
Cascade economics (ex9) & all-frontier: $25.8$ per correct artifact,
$3.0\%$ defects & cascade: $7.6$ per correct artifact, $1.5\%$ defects \\
Composition chain (ex10) & model-only prevalence error grows $5\times$
after silent labeler drift & PPI coverage $89.1\%$ through the drift;
monitor alarms at median day 41 (onset 31) \\
\bottomrule
\end{tabular}
\end{table}

\begin{table}[t]
\centering\small
\caption{Perturbation checks: each demonstration's key assumption is
deliberately violated; the table reports what breaks and which alarm the
library fires. ``Shipped'' alarms are runnable diagnostics; the causal
row's alarm is procedural (a playbook-forced escalation, not a runtime
statistic). False-alarm and detection rates are Monte Carlo estimates over
200--300 replications.}
\label{tab:perturb}
\begin{tabular}{@{}L{3.2cm}L{3.6cm}L{3.4cm}L{3.2cm}@{}}
\toprule
Demonstration & Perturbation & Consequence & Alarm \\
\midrule
Verifier gate (ex1, ex1b) & post-calibration drift: new defect type scores
higher & defects among accepted $2.8\%\to12.0\%$; target breached in
$86\%$ of runs & shipped (ex1b): audit-stream monitor fires at median
day 31 (onset 21), false-alarm rate $5.7\%$ vs.\ $5\%$ target \\
Nested risk (ex5b) & inner-sample budget misallocated ($M{=}1$) &
$4.6\times$ overstatement of the risk probability at identical total
cost & shipped: budget-aware inner/outer allocation rule (bias
$\propto 1/M$ made explicit) \\
Causal effect (ex7) & confounder hidden from the analyst & DML/AIPW
biased: $4.54$ for a true $2.00$ --- identification, not machinery,
fails & procedural: playbook-forced identification statement; sensitivity
analysis or escalation \\
Queue summary (ex8) & batch length short vs.\ relaxation time & batch
means undercover ($83.3\%$ vs.\ $90\%$ nominal) & shipped: batch-length
vs.\ autocorrelation-time check \\
Composition (ex10) & upstream labeler silently degrades
(specificity $0.97\to0.85$) & model-only estimate error grows $5\times$,
silently & shipped: anytime-valid monitor on the rectifier stream fires
at median day 41, false-alarm rate $5.5\%$ vs.\ $5\%$ target \\
\bottomrule
\end{tabular}
\end{table}

\paragraph{Calibrated gates for LLM verifiers (deployed).} Any pipeline
that generates artifacts in bulk --- code patches, extracted records,
auto-generated assessment items --- ends in the same bottleneck: an LLM
judge accepts or rejects, and the judge's error rate is unknown and drifts
across domains. The \skillname{conformal-uq} gate calibrates the acceptance
threshold on a modest labeled sample and returns a finite-sample guarantee
of the form ``at most $5\%$ of auto-accepted artifacts are defective, with
$90\%$ confidence,'' with abstentions routed to a human
\citep{vovk2005,angelopoulos2023gentle}. The difference is not accuracy but
\emph{auditability}: the gate's claim survives a hostile review that ``the
judge seemed fine'' does not. This is runtime V\&V (Layer~3) for the most
common trust decision in agent systems --- trusting a second model's
judgment. \emph{Demonstration (ex1):} against a $5\%$ defect target, the
naive ``judge says pass'' gate ships $8.0\%$ defects and breaches the
target in every calibration draw; the calibrated gate ships $2.8\%$ and
breaches in $8.3\%$ of draws --- within its stated $10\%$ risk --- at the
price of accepting $51\%$ rather than $76\%$ of artifacts. The price is
visible, priced, and chosen. \emph{Guarantee class, alarm, loss:} exact
finite-sample (Clopper--Pearson risk-controlling threshold), valid under
exchangeability of calibration and deployment artifacts with a fixed
scoring rule. Because Table~\ref{tab:perturb} shows exactly that
assumption voiding the guarantee under drift ($12\%$ defects shipped,
silently), the gate does not deploy alone: its companion monitor (ex1b)
feeds a small daily gold audit of \emph{accepted} artifacts into an
anytime-valid boundary and converts post-calibration drift into an alarm
at a budgeted false-alarm rate ($5.7\%$ realized against a $5\%$ target,
median ten days to detection). The bound feeds a defect-cost versus
abstention-cost model: with defect cost $C$, auto-acceptance is worth it
precisely when $0.05\,C$ beats the per-artifact cost of human review.

\paragraph{The always-on dashboard.} Modern metric cultures monitor
continuously and act on the first significant-looking excursion; under
peeking, a fixed-horizon test's false-alarm guarantee is void, and with
enough looks a spurious alarm is certain. The \skillname{anytime-valid}
skill replaces the test with a confidence sequence or e-process
\citep{howard2021,ramdas2023}: valid at every time simultaneously, so the
team may stop, continue, or extend for any reason without invalidating the
inference, with online FDR control across the stream of hypotheses a
dashboard culture generates. \emph{Demonstration (ex2):} under the null, a
year of daily peeking at $\alpha=0.05$ produces at least one false alarm
$47.4\%$ of the time; a Robbins normal-mixture confidence sequence,
$2.2\%$. Under a real drift of $0.25\sigma$ per day (a deliberately large daily
standardized effect, chosen for a legible stopping-time distribution) the
same sequence detects in $97\%$ of runs at a median of day 117 --- eight
months earlier than the fixed-horizon design it replaces. \emph{Guarantee
class, alarm, loss:} exact time-uniform under the martingale assumption;
the crossing itself is the alarm; the loss trade is the false-alarm budget
against detection delay, set by $\alpha$ and the mixture parameter.

\paragraph{Measurement with model-generated labels.} An analyst wants the
prevalence of some property in fifty thousand documents --- churn-risk
support tickets, a defect category in generated content, a competency
signal in HR text --- and an LLM can label all of them for cents. Treating
those labels as truth imports a bias of unknown sign and size; using only
the few hundred human labels wastes the model.
\skillname{prediction-powered-inference} \citep{angelopoulos2023ppi}
returns a confidence interval centered by the cheap predictions and
\emph{corrected} by the gold set: near-model-scale precision, human-scale
validity. For LLM-heavy organizations this is arguably the single
highest-leverage skill, because it legalizes a workflow everyone is
already doing invalidly. \emph{Demonstration (ex3):} with true prevalence
$12\%$ and a sensitivity-$0.95$/specificity-$0.98$ labeler, the model-only
estimate concentrates tightly on $13.2\%$ --- its $90\%$ interval covers
the truth in $0\%$ of repetitions --- while the 300-label gold-only
interval is honest but 6.2 points wide. PPI delivers a 2.9-point interval
with $89.7\%$ coverage: the 300 human labels do the work of $\sim$1{,}350.
\emph{Guarantee class, alarm, loss:} asymptotic, requiring the gold sample
to be a random draw from the measured population; the rectifier's mean is
itself the drift sensor (ex10 monitors exactly this stream); the interval
feeds the loss behind the prevalence decision --- overstaffing cost
against missed-case cost at a triage or escalation threshold.

\paragraph{Designing the expensive experiment.} When measurements cost real
money --- sensor placement on an infrastructure network, the next assay in
a lab campaign, the next item administered in an adaptive test ---
\skillname{bayesian-optimal-design} ranks candidate designs by expected
information gain per unit cost, choosing the estimator from its ladder
according to the problem's structure \citep{long2013,beck2018,ryan2016},
with sequential campaigns handled as belief-state planning. The delivered
difference is fewer experiments to a target posterior precision; the
delivered \emph{discipline} is that the agent must state what the
experiment is for before spending on it. \emph{Demonstration (ex4):}
calibrating a three-parameter response curve to a target posterior
precision takes the greedy max-EIG design 22 measurements against a median
of 37 for random placement --- a $40\%$ saving --- while the intuitive
``measure the center repeatedly'' design \emph{never} gets there: its
information matrix is singular in the slope and curvature directions, a
failure EIG accounting makes visible before the budget is spent.

\paragraph{Certifying against rare failures.} A reliability engineer must
certify that a failure probability --- a grid contingency, a network
outage, a safety-case event --- is below $10^{-6}$; naive Monte Carlo needs
on the order of $10^{8}$ simulations to even see the event. The
\skillname{rare-event-is} machinery (control-based importance sampling,
splitting for dynamic systems) reduces the required simulation budget by
orders of magnitude, while its efficiency diagnostics alarm when the
proposal itself is unreliable \citep{bucklew2004,rubino2009}. Downstream of
the tail probability sits the decision-relevant quantity, and
\skillname{risk-measure-estimation} owns it: VaR/CVaR and expected
shortfall estimated by nested and multilevel schemes
\citep{rockafellar2000,giles2019}, with elicitability theory
\citep{gneiting2011} supplying the backtests. Together they are the loss
side of the base contract: not ``how likely is failure,'' but ``what should
we provision for it.'' \emph{Demonstrations (ex5, ex5b):} for a Gaussian
tail probability of $1.02\times 10^{-6}$, naive Monte Carlo with a million
samples sees a single hit ($99\%$ relative standard error); a tilted
importance sampler at the same budget achieves $0.23\%$ --- a variance
reduction of $1.8\times 10^{5}$, or $9.8\times 10^{7}$ versus $540$ samples
for $10\%$ relative error. On the nested side, estimating a $5\%$
probability of large conditional loss with a single inner sample returns
$0.231$ --- a $4.6\times$ overstatement invisible to the user --- while a
budget-aware inner/outer allocation at \emph{identical total cost} returns
$0.053$ with $10\times$ lower RMSE.

\paragraph{Reporting the winners.} An agent screens two thousand features,
five hundred generated strategies, or a hundred subgroups, and reports the
top ten with confidence intervals --- the winner's curse, mechanized at
agent speed. \skillname{selective-inference} makes the shortlist honest:
FDR control over the screen \citep{bh1995,barber2015,candes2018} and
selection-adjusted intervals on what survives. The organizational symptom
this cures is familiar to every analytics team: effects that shrink or
vanish on replication. \emph{Demonstration (ex6):} screening 2{,}000 null
effects at $|z|>2.58$ and reporting naive $90\%$ intervals on the survivors
yields \emph{zero} coverage --- any screen stricter than the interval
half-width guarantees a miss, which is the winner's curse in its purest
mechanical form --- while conditional truncated-normal intervals restore
$90.4\%$. On a mixed panel, a Benjamini--Hochberg shortlist certifies its
false discovery rate ($9.2\%$ realized against a $10\%$ budget) and finds
\emph{more} true signals than the uncertified naive top-$k$ it replaces.
\emph{Guarantee class, alarm, loss:} exact conditional coverage and
finite-sample FDR control; the loss trade is follow-up cost per false
discovery against the value of a missed true signal, which sets the FDR
budget and the shortlist size.

\paragraph{Targeting instead of averaging.} Retention offers, training
interventions, and pricing changes all pose the same question: not ``does
it work on average'' but ``for whom.'' \skillname{causal-ml} carries the
modern answer --- double/debiased ML for confounding-robust effects
\citep{chernozhukov2018}, causal forests for heterogeneity
\citep{wager2018}, and policy learning with valid inference
\citep{athey2021} --- so that treatment is assigned where the effect is,
under stated assumptions the playbook forces the agent to check rather than
assume. \emph{Demonstration (ex7):} with treatment assignment confounded by
a single observed covariate, the difference in means reads $6.23$ for a
true effect of $2.00$ --- a $3\times$ overstatement that would survive any
code review --- while cross-fitted AIPW returns $2.00\pm0.04$. In the
targeting variant, treating only the subgroup whose estimated effect
exceeds cost doubles the realized policy value over treat-everyone
($+1.00$ versus $+0.50$ per unit). \emph{Guarantee class, alarm, loss:}
asymptotic/semiparametric, and conditional on identification --- the
playbook forces a written statement of the assignment mechanism, observed
confounders, and overlap diagnostics, and Table~\ref{tab:perturb} shows
why: with a hidden confounder the same machinery returns $4.54$ for a true
$2.00$, so when unmeasured confounding is plausible the skill's output is
a sensitivity analysis or an escalation, not an estimate. The loss hookup
is explicit: treatment cost against estimated subgroup effect.

\paragraph{Trusting a simulation's summary.} A discrete-event or
digital-twin simulation drives a capacity decision. The model may be
perfectly correct and the summary still invalid: autocorrelated output
averaged as if independent, the initial transient left in, tails read from
a single trajectory. \skillname{simulation-output-analysis} supplies the
output-side discipline --- batch-means and regenerative confidence
intervals, transient handling, and Rhee--Glynn randomization
\citep{asmussen2007,rhee2015} for unbiased averaging across parallel
replications --- while \skillname{measure-transport-inference} keeps the
twin \emph{calibrated}, updating its posterior against field data without
an MCMC campaign \citep{elmoselhy2012,marzouk2016}.
\emph{Demonstration (ex8):} for an M/M/1 queue at $90\%$ utilization (true
mean wait exactly 9), the iid-style $90\%$ interval from a
40{,}000-customer run covers the truth $5.0\%$ of the time --- and is
seductively narrow ($0.16$ wide) --- while batch means covers $83.3\%$ at
an honest width of $2.7$. The residual gap below nominal is itself
informative: batch length relative to the queue's relaxation time is the
diagnostic the skill checks before trusting the interval.
\emph{Guarantee class, alarm, loss:} asymptotic ($t$-interval on batch
means) with the batch-length diagnostic as the shipped alarm; the interval
feeds overprovisioning cost against shortfall or SLA-violation cost in the
capacity decision.

\paragraph{Paying for intelligence correctly (deployed).} The economics of
LLM systems are themselves a V\&V surface. \skillname{model-hierarchy}
reframes the managed quantity from cost per call to cost per
\emph{verified} artifact: bulk generation at the cheapest level that
verification can police, escalation only on disagreement or failed checks
--- an organizational analogue of multilevel Monte Carlo's telescoping
principle \citep{giles2008,giles2015} with disagreement as the
a-posteriori error indicator. The library's own construction
(Section~\ref{sec:harvest}) is the existence proof; the pattern applies to
any bulk LLM workload where checking is cheaper than producing.
\emph{Demonstration (ex9):} with a cheap model at $88\%$ accuracy and unit
cost, a frontier model at $97\%$ and $25\times$ cost, and a verifier with
$92\%$ error recall at $2\times$ cost, the verify-and-escalate cascade
dominates the all-frontier strategy on \emph{both} axes: $1.5\%$ residual
defects (versus $3.0\%$) at $7.6$ per correct artifact (versus $25.8$).
Once a plausible downstream cost per shipped defect is added, the
all-cheap strategy --- nominally $25\times$ cheaper --- becomes more
expensive than the cascade as well.

\paragraph{Composition: a three-skill chain under drift.} Skills are used
in pipelines, and validity does not compose automatically --- which is why
the library's composition discipline is demonstrated rather than assumed.
The chained demonstration (ex10) is the workload every LLM-heavy
organization runs: a labeler processes two thousand documents a day, a
40-document gold audit is drawn daily, and the pipeline must deliver a
daily prevalence estimate \emph{and} notice when the labeler silently
changes. Link one is \skillname{prediction-powered-inference}, whose daily
interval is valid regardless of labeler quality because the rectifier is
estimated on the same day's gold pairs. Link two is an
\skillname{anytime-valid} monitor on the rectifier stream, calibrated on a
ten-day burn-in, whose boundary runs in an intrinsic time that charges for
the estimated baseline's uncertainty. At day 31 the labeler's specificity
silently drops ($0.97\to0.85$). \emph{Measured:} the model-only estimate's
error grows fivefold with no visible symptom; the PPI intervals hold
$89.1\%$ coverage straight through the drift; and the monitor fires at a
median of day 41 with a $5.5\%$ false-alarm rate against its $5\%$ target
--- the upstream assumption violation converted into an explicit,
budgeted alarm. This is the framework's composition claim in full: chains
are safe exactly insofar as each link's preconditions are watched by
another link, and contract-level composition semantics beyond this
pattern are an open frontier (Section~\ref{sec:scope}).

\section{Building the library: curriculum harvesting}\label{sec:harvest}

The toolkit above was not assembled from skill marketplaces --- which are
dominated by software-engineering workflows and rarely ship falsifiable
claims --- but harvested from the literature. The unit of harvesting is the
\emph{methodological school}: an author or small cluster whose students and
collaborators propagate a coherent, transferable toolkit (multilevel Monte
Carlo around Giles and Tempone; measure transport around Marzouk; conformal
prediction and selective inference around Cand\`es; anytime-valid inference
around Ramdas). A school, unlike a paper, comes with an implicit
\emph{curriculum} --- the follow-up papers, the surveys, the software
habits, the standard counterexamples a practitioner needs to apply the
method without a supervisor --- and the harvesting protocol is built to
extract exactly that. Five stages; one school (occasionally a deliberately
paired pair of lineages) per session block.

\subsection{Stage 1: the author roster}

A live, git-tracked queue of candidate schools with explicit priorities.
Admission requires four criteria: \textbf{school-forming} (students and
collaborators propagate a coherent toolkit); \textbf{method-first}
(techniques transfer beyond the home domain); \textbf{practitioner-reachable}
(papers state algorithms, not only theorems); and \textbf{gate-plausible}
(at least one cluster could plausibly be admitted against the current
library --- checked before a row is added). Nominations come from the
operator's field knowledge, from coauthor-graph exploration of
already-harvested schools (resolving authors by OpenAlex identifiers from
joint works, never by name search --- name collisions and merged author
records are both incidents in the project log), and from cross-family
elicitation runs in which several cheap models from different providers
independently nominate schools against the same criteria and exclusion
list, with same-family agreement discounted as correlated evidence.

\subsection{Stage 2: scoped harvest into a citation-disciplined wiki}

A run opens with a locked scope (the clusters to harvest, named in the
roster row) and fans out parallel subagents over the school's literature.
The output is a \emph{wiki}: per-cluster pages with per-claim citations,
compiled under the rule that a claim without a source does not enter, and
checked by an independent verifier pass that samples claims and attempts to
falsify them against the cited sources. The wiki --- typically 20--35
admitted sources per school --- persists as the audit trail for every
downstream decision.

\subsection{Stage 3: the novelty gate}

Each method cluster receives a structured verdict against the existing
library: \textsc{adopt} (new skill, with a recorded ownership boundary
against every adjacent skill), \textsc{merge} (versioned delta to an
existing skill), or \textsc{defer} (not built, but with recorded trigger
conditions under which the verdict reopens --- deferral with triggers keeps
rejection cheap to reverse). The gate is adversarial by design: its default
answer is no, and overlap is scored before value. In the runs reported here
it rejected, among others, a standalone probabilistic-graphical-models
skill (thin wedge over existing Bayesian machinery) and a
wireless-communications performance-analysis school (a domain school whose
methodological content already reached the library through a method-first
lineage).

\subsection{Stage 4: build with standard-library scripts}

Winning clusters become skills: a playbook stating when the method applies,
the decision points, the standard failure modes, and the stopping rules,
plus executable scripts kept dependency-light (Python standard library
wherever feasible). The discipline is not aesthetic: skills travel across
machines and sandboxes, and a script that fails at import time spends the
agent's error budget before the statistics begin.

\subsection{Stage 5: verification anchors}

Before admission, each skill's scripts must reproduce a known ground truth
(Layer~1 of the framework); Section~\ref{sec:verify} gives the anchors
used. Admission control plus anchors is what separates harvesting from
``ask the model to summarize the literature'': every claim is cited, every
adoption is adjudicated, and every shipped procedure is falsifiable.

\section{Schools harvested and skills derived}\label{sec:schools}

Table~\ref{tab:schoolruns} summarizes the eight harvesting runs (July
2026): eleven new skills and six versioned upgrades. Per-school detail ---
run narratives, gate records, and the coauthor-graph exploration --- is
companion-paper material; here we note only what bears on the V\&V
program.

\begin{table}[t]
\centering\small
\caption{Harvesting runs (July 2026). ``Yield'' lists new skills and
version-bumped merges.}
\label{tab:schoolruns}
\begin{tabular}{@{}L{3.4cm}L{4.2cm}L{6.6cm}@{}}
\toprule
School & Clusters harvested & Yield \\
\midrule
R.~Tempone (KAUST/RWTH) & MIMC; optimal experimental design; control-variate
importance sampling; adaptivity & \skillname{bayesian-optimal-design} (new);
\skillname{rare-event-is} (new); \skillname{model-hierarchy} v1.2
(multi-index); \skillname{numerical-vv} v1.1 (adaptivity checks) \\
\addlinespace
Y.~Marzouk (MIT) & Transport maps; likelihood-informed subspaces; sequential
OED; transport filtering & \skillname{measure-transport-inference} (new);
\skillname{bayesian-optimal-design} v1.1 (sequential OED);
\skillname{bayesian-workflow} v1.1 (high-dim.); filtering \textsc{defer}red
with triggers \\
\addlinespace
E.~Cand\`es (Stanford) & Conformal prediction; knockoffs; selective
inference; LLM verifier gates & \skillname{conformal-uq} (new);
\skillname{selective-inference} (new) \\
\addlinespace
M.~I.~Jordan (Berkeley) & Variational inference; PGMs; decision-focused ML &
\skillname{prediction-powered-inference} (new); \skillname{bayesian-workflow}
v1.2 (VI); PGMs rejected standalone \\
\addlinespace
Athey--Imbens--Wager + Chernozhukov (Stanford/MIT) & Causal forests; DML;
policy learning; adaptive experiments & \skillname{causal-ml} (new);
\skillname{study-design} v1.1 (adaptive-experiment checklist) \\
\addlinespace
A.~Ramdas (CMU) & Confidence sequences; e-processes; online FDR &
\skillname{anytime-valid} (new); \skillname{selective-inference} v1.0.1
(boundary) \\
\addlinespace
Glynn + Rubino (Stanford/INRIA) & Output analysis; unbiased multilevel
randomization; splitting; heavy tails & \skillname{simulation-output-analysis}
(new); \skillname{rare-event-is} v1.1 (splitting, heavy tails, efficiency
gates) \\
\addlinespace
Haji-Ali + Krumscheid (Heriot-Watt/KIT; Giles bridge) & Nested MC for risk
functionals; MLMC for distributions; robust UQ &
\skillname{risk-measure-estimation} (new) \\
\bottomrule
\end{tabular}
\end{table}

Three features of the runs bear directly on the framework's credibility.
First, the gate rejected and merged as often as it adopted: variational
inference entered as a version bump to an existing Bayesian-workflow
skill; probabilistic graphical models were rejected standalone as a thin
wedge; sequential experimental design, left open by one school's run, was
closed by another school's belief-space formulation as a merge rather than
a new skill --- admission control functioning as designed. Second, the
Cand\`es-school run is where the library became self-policing: conformal
prediction is itself a V\&V technology (distribution-free finite-sample
coverage from exchangeability alone \citep{vovk2005,lei2018}), and the
calibrated verifier gates it contributed now gate the library's own
construction pipeline. Third, the roster is not closed under operator
knowledge: the risk-measure school was surfaced by coauthor-graph
exploration, filling a gap the base contract had made visible --- the
library could estimate probabilities and expectations but did not own
coherent-risk-measure estimation \citep{rockafellar2000,giles2019}, the
loss side of ``every decision names its loss.''

\section{Verification discipline in practice}\label{sec:verify}

The anchor rule --- every skill that ships numeric code must reproduce a
ground truth exactly known by other means, on the same code path a user
will exercise --- is Layer~1 of the framework. Representative anchors:

\begin{itemize}[leftmargin=1.6em]
\item \textbf{Closed-form EIG} (\skillname{bayesian-optimal-design}): for
linear-Gaussian models the expected information gain has a closed form; the
Monte Carlo estimator ladder must converge to it, and the anchor doubles as
a demonstration of the inner-loop bias that makes naive nested Monte Carlo
drift with inner sample size \citep{ryan2016}.
\item \textbf{Cholesky as transport map}
(\skillname{measure-transport-inference}): for Gaussian targets the
Knothe--Rosenblatt map is triangular and linear; the fitted map must
recover the Cholesky factor.
\item \textbf{Exact finite-sample conformal coverage}
(\skillname{conformal-uq}): split conformal at level $\alpha$ with $n$
calibration points has coverage exactly
$\lceil (n{+}1)(1{-}\alpha)\rceil/(n{+}1)$ under exchangeability, with a
Beta-distributed conditional coverage \citep{vovk2005}. This anchor is
sharp enough to catch off-by-one errors in the quantile index --- and did:
the first generated implementation passed a loose ``coverage
$\approx 1-\alpha$'' check but failed the exact band, revealing precisely
the \mbox{$\lceil\cdot\rceil$-vs-index} bug the anchor was designed to
catch.
\item \textbf{Winner's curse, quantified} (\skillname{selective-inference}):
on selected null effects, naive 90\% confidence intervals achieved 20.6\%
coverage in the skill's bundled demonstration, which selects by top-$k$
rank; the paper's ex6 uses a threshold screen instead, under which naive
coverage is exactly zero. The two figures are different selection rules,
both shipped --- the skill carries the failure, not only the fix.
\item \textbf{Variance-reduction accounting} (\skillname{rare-event-is}):
the skill's bundled control-based importance-sampling demonstration
documents a variance reduction of order $10^4$ on its test problem (the
paper's ex5, a different and easier tail, reaches $1.8\times10^{5}$)
alongside the structural fact the anchor verifies: an imperfect control
changes \emph{variance}, never \emph{bias} --- the property that makes the
method safe to hand to a non-expert.
\end{itemize}

Two lessons generalize. First, anchors should be \emph{exact}, not
approximate: the conformal bug survived an approximate check and died on
the exact one --- the direct analogue of the observation in classical V\&V
that qualitative ``looks right'' comparisons pass codes that
order-of-convergence tests kill \citep{roache1998}. Second, skills should
ship their canonical \emph{failure} as a runnable demonstration alongside
the fix; for statistical methods, the counterexample is half the pedagogy.

\section{Orchestration and economics}\label{sec:orchestration}

The construction pipeline runs on a hierarchy of models orchestrated by the
\skillname{model-hierarchy} skill, in deliberate analogy with multilevel
Monte Carlo \citep{giles2008,giles2015}: expensive evaluations are reserved
for what only they can do, and the bulk of the sampling happens at the
cheap levels. Tier-0 leaves (free or near-free models across several
providers, plus local models) perform bulk generation --- literature
harvesting, candidate nomination, replicate drafting --- with replication
across model \emph{families} and same-family agreement discounted as
correlated evidence; a mid tier verifies, dedupes, and aggregates; a root
model locks scope, arbitrates the novelty gate, and owns the audit trail.
A strict generator/verifier separation is enforced: no tier adjudicates its
own output.

Two governance layers make this safe to run unattended. A
\emph{data-egress gate} classifies every payload before it leaves the
machine (free-tier providers may log prompts; only public,
published-literature content is dispatched to them). And an append-only run
log with per-call cost accounting keeps the economics honest: an external
multi-model review panel over the library's entire design cost \$0.61
across three rounds --- the marginal cost of adversarial review at this
scale is effectively zero, which changes how often one should convene it.

The human operator's role concentrates at three points: nominating and
signing off roster rows, resolving gate debates the adversarial review
leaves open, and supplying the loss function --- deciding what the library
is \emph{for}, which no amount of literature harvesting determines.
Everything between those points is delegated, logged, and reversible.

\paragraph{Self-application.} A paper advocating falsifiable checks owes a
brief statement of which checks were applied to its own artifacts --- with
the caveat that the manuscript is a human-authored argument, not a runtime
artifact, and was not produced by the pipeline it describes. What was
checked: the library's skill scripts passed their anchors before
admission; the demonstration scripts are anchor-disciplined (fixed seeds,
exact ground truths) and ship with the paper; and the manuscript underwent
a documented three-round adversarial review by seven frontier models from
seven providers (costs and verdicts archived in the project's run logs),
whose surviving criticisms produced, among other revisions, the renaming
of Layer~2, the scope section, the perturbation table, the gate's
companion monitor, and the composition demonstration. The raw review
materials remain in the private project archive and are not part of the
public repository. What has not
happened: independent replication by a second operator, and validation
against a live workload --- both named as the natural next studies. Until
then this remains what it claims to be: a case report whose load-bearing
claims are third-party falsifiable from the public artifacts.

\section{The mechanism census}\label{sec:census}

The most useful strategic artifact to emerge from the program is a one-line
characterization of each school by the structural resource its methods
\emph{spend} to buy statistical efficiency or validity
(Table~\ref{tab:census}). The census answers the question that otherwise
dooms library growth to redundancy: a new school is worth harvesting only
if it spends a resource the library does not yet spend. It functions both
as a diversification rule for the roster and as a stopping rule: when
nominations keep landing on spent mechanisms, the literature's marginal
value to the library has collapsed.

\begin{table}[t]
\centering\small
\caption{Mechanism census: the resource each harvested school spends.}
\label{tab:census}
\begin{tabular}{@{}L{4.6cm}L{9.6cm}@{}}
\toprule
School & Resource spent \\
\midrule
Tempone (MLMC/MIMC, OED) & \emph{Error indicators}: budget allocated along
estimated error contributions across levels and indices \\
Marzouk (transport) & \emph{Geometry}: low-dimensional and triangular
structure of the posterior \\
Cand\`es (conformal, knockoffs) & \emph{Exchangeability}: symmetry of the
data under permutation, spent for finite-sample validity \\
Ramdas (anytime-valid) & \emph{Martingale structure}: nonnegative
supermartingales spent for validity under optional stopping \\
Jordan (PPI, VI) & \emph{Cheap predictions}: abundant ML output debiased
against a small gold-standard set \\
Athey--Wager--Chernozhukov & \emph{Orthogonality and design}: Neyman
orthogonality and randomization spent for robustness to nuisance error \\
Glynn--Rubino & \emph{Probabilistic structure of the simulator}:
regeneration, stationarity, and coupling spent for valid output analysis \\
Haji-Ali--Krumscheid & \emph{Smoothness of the functional}: regularity of
risk functionals spent inside nested and multilevel estimators \\
\bottomrule
\end{tabular}
\end{table}

\section{Limitations}\label{sec:limitations}

This is a single-operator case report, and several limitations are
structural. (i)~\emph{No controlled evaluation}: we report that anchors
caught bugs and gates prevented duplication, but not a counterfactual
library grown without them; a matched comparison is the natural next
study. (ii)~\emph{Roster bias}: nomination begins from the operator's field
knowledge and the coauthor graphs of already-harvested schools, which
biases the library toward the Monte Carlo/UQ lineage it started from; the
cross-family elicitation runs are a partial correction, not a cure.
(iii)~\emph{Bibliometric fragility}: author-disambiguation incidents (name
collisions; merged OpenAlex author records) occurred and were caught only
because the protocol requires resolving authors by identifiers from joint
works and verifying works individually. (iv)~\emph{Staleness}: model
training knowledge dates affiliations and ``current'' research lines; the
roster convention is that all such fields are unverified until harvest
time. (v)~\emph{Anchor coverage}: an anchor verifies a code path, not a
skill's prose; a playbook can still misstate when a method applies.
Adversarial multi-model review mitigates but does not close this gap ---
in V\&V terms, Layer~1 is necessary but never sufficient for Layer~2.
(vi)~\emph{Stylized demonstrations}: the numerical examples of
Section~\ref{sec:toolkit} use simulated data-generating processes chosen
so that ground truth is exact and the contrast is legible; they demonstrate
the \emph{mechanism} of each failure and its repair, not the effect size a
particular applied dataset will exhibit. Real problems will sit anywhere
between the two columns of Table~\ref{tab:numbers}, and the honest claim is
only that the naive column's failure modes are generic, not that its
magnitudes are. (vii)~\emph{Cooperative-agent assumption}: the framework
addresses statistical failure in pipelines whose components are fallible
but not adversarial; prompt injection or an agent actively optimizing
against its own guardrails is a security problem outside this paper's
threat model.

\section{Related work}\label{sec:related}

The V\&V tradition this paper transplants is codified in
\citet{roache1998} and \citet{oberkampf2010} for computational science and
in \citet{sargent2013} for simulation models; our contribution is the
mapping of its layers onto LLM-assisted statistical work and a concrete,
skill-packaged implementation. The skill mechanism and its
progressive-disclosure design are described in \citet{anthropic2025skills};
tool learning and tool-use augmentation of language models are surveyed in
\citet{schick2023,mialon2023}. The library differs from
retrieval-augmented approaches \citep{lewis2020} in producing durable,
versioned, executable procedure rather than answer-time context, and from
automatic skill discovery in embodied agents \citep{wang2023voyager} in
sourcing procedures from a scientific literature under human admission
control. The anchor discipline is a relative of established
statistical-software testing practice --- notably simulation-based
calibration for Bayesian computation \citep{talts2018} and the tradition
of oracle-based testing of statistical routines --- applied here at the
level of an agent's entire procedural repertoire rather than a single
sampler. Several skills package methods that are themselves recent
responses to machine-learning-era validity problems: conformal prediction
\citep{vovk2005,angelopoulos2023gentle}, prediction-powered inference
\citep{angelopoulos2023ppi}, and anytime-valid inference
\citep{ramdas2023}; the library's role is to make them \emph{procedurally
available} to agents at the moment of use, with their assumptions checked
rather than assumed.

\section{Conclusion}\label{sec:conclusion}

The reliability problem of LLM-assisted analysis is not that models
hallucinate facts --- retrieval and citation discipline address that ---
but that pipelines emit statistically invalid numbers with perfect
confidence. The remedy proposed here is neither new models nor new
benchmarks but an old discipline with a strong track record in exactly this
situation: verification and validation, applied at three layers ---
anchor-verified tooling, operating-characteristic-verified procedures, and
runtime guarantee-or-alarm on every delivered number. The stochastics and UQ literature supplies both the
machinery and, through its exactly solvable special cases, the anchors that
make the machinery falsifiable; the agent-skill format supplies the carrier
that makes it durable, versioned, and auditable; and curriculum harvesting
from methodological schools supplies the economics that make the whole
program feasible for a single operator in days rather than quarters. The
measured gaps of Table~\ref{tab:numbers} are the argument in brief: the
naive column is documented default behavior of unassisted pipelines today,
and the distance to the other column is the loss available to anyone who
leaves those failure modes unguarded. The framework's own contribution is
not the statistics --- the methods are decades old and belong to their
schools --- but making the guarded column agent-discoverable,
runtime-audited, falsifiable by anyone, and alarmed at its known
boundaries (Table~\ref{tab:perturb}).

\medskip
\noindent A closing note on provenance: the coauthor-graph stage of the
harvesting program eventually surfaced the author's own doctoral work
\citep{moraes2014} among the lineages mined --- the toolchain whose V\&V
habits this library generalizes, already installed by construction.

\section*{Acknowledgments}

The skill library, the harvesting runs, the numerical demonstrations, and
the drafting of this manuscript were carried out with substantial
assistance from large language models (Anthropic's Claude family as root
and mid-tier models; various providers' models via OpenRouter as tier-0
leaves), under the orchestration and verification protocol described in
the paper. All verdicts, roster decisions, and claims were reviewed by the
human author, who takes full responsibility for the content.

\bibliographystyle{plainnat}
\bibliography{references}

\end{document}
